% One setpoint, from the host's loop to the node and back, with the two
% timestamps marked.
\begin{tikzpicture}[font=\sffamily\scriptsize, x=1cm, y=1cm]

\foreach \x/\name/\id in {0/{setpoint source}/src, 3.2/{the kernel}/krn,
                          6.4/{the adapter}/adp, 9.6/{the node}/nd,
                          12.8/{the control loop}/loop}{
  \node[actor, text width=21mm] (\id) at (\x,0) {\name};
  \draw[lifeline] (\id.south) -- ++(0,-8.4);
}

\draw[msg] (0,-1.0) -- node[above]{send on a raw socket} (3.2,-1.0);
\node[activation, minimum height=3mm] at (3.2,-1.2) {};
\draw[msg] (3.2,-1.9) -- node[above]{out through the network device} (6.4,-1.9);
\draw[msg] (6.4,-2.6) -- node[above]{on the wire} (9.6,-2.6);
\node[activation, minimum height=4mm] at (9.6,-2.85) {};

\node[tlabel, anchor=west, text=blue!55!black] at (9.75,-3.4)
  {timestamp: hardware, at the start-of-frame bit};

\draw[msg] (9.6,-4.0) -- node[above]{the nearest period takes it} (12.8,-4.0);

\draw[reply] (12.8,-4.9) -- node[below]{state frame, every period} (9.6,-4.9);
\draw[msg] (9.6,-5.5) -- node[above]{on the wire} (6.4,-5.5);
\draw[msg] (6.4,-6.1) -- node[above]{interrupt, then a transfer if it is on USB} (3.2,-6.1);
\node[activation, minimum height=3mm] at (3.2,-6.3) {};

\node[tlabel, anchor=east, text=red!60!black] at (3.05,-6.8)
  {timestamp: the kernel, when the interrupt is served};

\draw[msg] (3.2,-7.4) -- node[above]{recv} (0,-7.4);

\node[umlnote, text width=52mm, anchor=north west] at (0,-9.0)
  {The two timestamps are not equivalent and the chapter measures the
   difference rather than assuming it. The node's is taken by hardware before
   anything else happens. The host's is taken after the adapter has delivered
   the frame, which on a USB adapter includes a transfer that was scheduled
   rather than immediate.};
\node[umlnote, text width=52mm, anchor=north west] at (6.6,-9.0)
  {The host's loop is not a real-time loop and does not pretend to be. It
   produces setpoints at a rate the node can absorb, and the node's own control
   period keeps time. A host that tries to be the clock puts its own scheduling
   jitter into the joint's motion.};
\end{tikzpicture}
